\makesectionframe{NLP}[Natural Language Processing]

\begin{contentframe}
    \frametitle{NLP}

    \begin{itemize}
        \item Molte librerie con tante funzionalità
        \begin{itemize}
            \item \texttt{nltk, spacy, textblob, ...}
        \end{itemize}

        \bigskip
        \item Alcune funzionalità sono fornite da più librerie
    \end{itemize}
\end{contentframe}

\begin{exampleframe}
    \frametitle{NLP}
    \framesubtitle{Note}

    \begin{itemize}
        \item Spesso è necessario scaricare \textit{corpora} o altri dati
        \item I file sono salvati sul computer, quindi non serve scaricarli ogni volta
    \end{itemize}
\end{exampleframe}

\begin{exerciseframe}
    \frametitle{Installazione}

    \begin{enumerate}
        \item Aggiungere nel file \texttt{requirements.txt} le seguenti librerie:
        \begin{itemize}
            \item \texttt{nltk, spacy, textblob}
        \end{itemize}
        \item Installarle nell'ambiente virtuale
        \item Eseguire il seguente codice nel terminale (dopo aver attivato l'ambiente):
        \begin{itemize}
            \item \texttt{python -m spacy download en\_core\_web\_sm}
        \end{itemize}
    \end{enumerate}
\end{exerciseframe}

\begin{exerciseframe}
    \frametitle{Progettino}

    \begin{enumerate}
        \item Nella classe \texttt{Text}:
        \begin{itemize}
            \item Modificare \texttt{self.words}, usando \texttt{nltk.tokenize.word\_tokenize(...)}
            \pause
            \item Aggiungere \texttt{self.sentences}, che tramite \texttt{nltk} divide correttamente in frasi
            \item Aggiungere \texttt{get\_random\_sentence(self)}, che ritorna una frase casuale nel testo
            \item Aggiungere \texttt{get\_random\_verb(self)}, che ritorna un \textbf{verbo} casuale nel testo
        \end{itemize}
    \end{enumerate}
\end{exerciseframe}

\begin{exerciseframe}
    \frametitle{Progettino}

    \begin{enumerate}
        \setcounter{enumi}{1}
        \item Creare il file \texttt{main.py}, in cui:
        \begin{itemize}
            \item Si importa il file \texttt{texts.py} 
            \item Si analizzano alcuni testi \textit{(esempio guidato)}
            \begin{itemize}
                \item Named Entities
                \item Analisi grammaticale
                \item Sentiment Analysis
                \item Stemming / Lemmatization
            \end{itemize}
        \end{itemize}
    \end{enumerate}
\end{exerciseframe}

% caricare cartella su teams
% slide: elencare brevemente le principali librerie 
% mostrare esempi interattivi
% per ogni altra cosa, ricecare su internet
% esempio: usare T5 per riassumere testi


\makesectionframe{API ChatGPT}

\begin{contentframe}
    \frametitle{API ChatGPT}

    \begin{itemize}
        \item Si può automatizzare l'interazione con ChatGPT
        \begin{itemize}
            \item Effettuare richieste in automatico
            \item Effettuare automaticamente azioni sulle risposte
        \end{itemize}

        \bigskip
        \item Libreria necessaria: \texttt{openai} 
        \item Requisito: chiave API di ChatGPT \textit{(a pagamento)}
    \end{itemize}
\end{contentframe}

\begin{contentframe}
    \frametitle{Struttura dei prompt}

    \begin{columns}
        \col{.5}
        \begin{itemize}
            \item Lo strumento genera il messaggio a partire dai precedenti
    
            \bigskip
            \item Tre tipi di messaggio
            \begin{itemize}
                \item \textbf{Sistema:} istruzioni per il modello \textit{(utile solo in situazioni particolari)}
                \item \textbf{Utente:} richiesta dell'utente
                \item \textbf{Assistente:} risposta del modello
            \end{itemize}
        \end{itemize}

        \col{.5}
        \centering
        \smartartflowvertical[System,User,Assistant]
    \end{columns}
\end{contentframe}

\begin{exampleframe}
    \frametitle{Struttura dei prompt}
    \framesubtitle{Esempio}

    \begin{itemize}
        \item \href{https://chatgpt.com/share/6759a21f-19ec-8003-a172-e0630eae761d}{Prompt 1} ( $\boxed{S + U} \rightarrow A$ )
        \item \href{https://chatgpt.com/share/67605a8c-2fec-8003-8e82-c54f49687c1b}{Prompt 2} ( $\boxed{S + U_1 + A_1 + U_2} \rightarrow A_2$ )

        \bigskip
        \item Il prompt di sistema non è specificato, quindi è:\\
            \texttt{You are ChatGPT, a large language model trained by OpenAI. Answer as concisely and informatively as possible. Your knowledge is up to date until [cutoff date].}
    \end{itemize}
\end{exampleframe}

\begin{contentframe}
    \frametitle{Note}

    \begin{itemize}
        \item Sono presenti tanti parametri che influenzano lo stile di risposta
        \begin{itemize}
            \item Temperatura, penalità di frequenza, top\_p, lunghezza massima, ...
        \end{itemize}
    
        \bigskip
        \item La conversazione è semplicemente una sequenza di stringhe
        \begin{itemize}
            \item Non devo necessariamente generarla tutta tramite AI
        \end{itemize}

        \bigskip
        \item I costi sono basati sul numero di token ($out > in$) e sul modello usato
    \end{itemize}
\end{contentframe}

\begin{exampleframe}
    \frametitle{API ChatGPT}

    \begin{itemize}
        \item Documentazione: \url{https://platform.openai.com/docs/api-reference/chat}
    \end{itemize}
\end{exampleframe}

\begin{exerciseframe}
    \frametitle{Progettino}
    \framesubtitle{Esercizio guidato}

    \begin{enumerate}
        \item Creare il file \texttt{chatgpt.py}, come da esempio guidato
        \item Creare la classe \texttt{AIGeneratedText}, come da esempio guidato
        \item Effettuare le stesse analisi di prima su testi generati da AI
    \end{enumerate}
\end{exerciseframe}

